% ============================================================
% CHAPTER 9: FRONTEND SYSTEM ARCHITECTURE AND IMPLEMENTATION
% ============================================================

The AgriVerify frontend provides a responsive and secure browser interface that connects users to the backend services. It is designed to accommodate two distinct user profiles: farmers accessing the system via mid-range Android devices over unstable field networks, and state administrators managing comprehensive supply chain dashboards. The technology stack Next.js~16, React~19, TypeScript~5 and Tailwind CSS was selected to ensure optimal client-side performance, accessibility and clean architectural separation of concerns without sacrificing maintainability \cite{salim2024microservices, singh2026impact}.

% ============================================================
\section{Frontend Architecture and Design Patterns}
\label{sec:frontend-architecture}

The architecture separates UI rendering from data fetching, API communication and authorization. Each concern lives in its own layer.

\subsection{Architectural Layers}
A request passes through four layers before reaching backend services (Figure~\ref{fig:frontend_architecture}):
\begin{enumerate}
    \item \textbf{Presentation UI (React \& Next.js):} Handles component rendering, form validation and responsive layout \cite{react2024}.
    \item \textbf{Client-Side State (TanStack Query):} Manages async data fetching, caching and optimistic mutations \cite{tanstack2024}.
    \item \textbf{Service Client (Axios):} Wraps network calls, standardising payloads and error handling.
    \item \textbf{API Proxy (Next.js Routes):} Routes client requests to ASP.NET Core endpoints without exposing internal URLs to the browser.
\end{enumerate}

\begin{figure}[H]
\centering
\begin{adjustbox}{max width=0.95\textwidth}
\begin{tikzpicture}[
    node distance=0.6cm and 1.2cm,
    every node/.style={
        draw=black!80, thick, rounded corners=3pt, align=center,
        minimum width=4cm, minimum height=0.9cm, font=\scriptsize, fill=white
    },
    arrow/.style={->, thick, >=stealth, draw=blue!80!black}
]
\node[fill=blue!10] (ui) {\textbf{UI Component Layer}\\(Next.js \& React)};
\node[fill=green!10, below=of ui] (hooks) {\textbf{State Synchronization}\\(TanStack Query)};
\node[fill=yellow!10, below=of hooks] (services) {\textbf{Service Layer}\\(Axios API Client)};
\node[fill=orange!10, below=of services] (proxy) {\textbf{API Proxy Route}\\(Next.js Intermediary)};
\node[fill=red!10, right=1.2cm of proxy] (backend) {\textbf{ASP.NET Core Backend}\\(Business \& Data Tier)};

\draw[arrow] (ui) -- (hooks);
\draw[arrow] (hooks) -- (services);
\draw[arrow] (services) -- (proxy);
\draw[arrow] (proxy) -- (backend);
\end{tikzpicture}
\end{adjustbox}
\caption{Frontend Layered Architecture and Request Execution Flow}
\label{fig:frontend_architecture}
\end{figure}

\subsection{Secure Proxy Architecture}
All outbound requests route through internal Next.js serverless endpoints (\texttt{/api/proxy/*}), which avoids CORS issues by keeping cross-origin communication server-side. Internal microservice addresses stay out of the browser, secrets remain on the server, and security headers are applied before responses reach the client \cite{masue2024quantifying}.

% ============================================================
\section{Technology Stack Analysis}
\label{sec:technology-stack}

Table~\ref{tab:frontend-stack} summarises the stack and the reasoning behind each choice.

\begin{table}[htbp]
\centering
\caption{Frontend Technology Stack}
\label{tab:frontend-stack}
\small
\renewcommand{\arraystretch}{1.1}
\begin{tabularx}{\textwidth}{|p{3.8cm}|X|}
\hline
\textbf{Technology} & \textbf{Architectural Reasoning} \\
\hline
Next.js 14 & Provides Server-Side Rendering and built-in API proxy routing~\cite{nextjs2024}. \\
\hline
React 18 & Declarative UI rendering with efficient Virtual DOM reconciliation. \\
\hline
TypeScript 5 & Compile-time type checking with backend DTO parity. \\
\hline
Tailwind CSS & Utility-first styling with minimal production bundle output~\cite{tailwind2024}. \\
\hline
Axios \& TanStack & HTTP pipeline management, token injection, and cache synchronization~\cite{tanstack2024}. \\
\hline
Zustand 4.5 & Lightweight atomic state slices for session persistence. \\
\hline
\end{tabularx}
\end{table}

% ============================================================
\newpage
\section{Modular Project Structure}
\label{sec:frontend-project-structure}

The source tree follows domain-driven modularity. Key directories: \texttt{src/app} (Next.js routing), \texttt{src/components} (reusable UI primitives), \texttt{src/hooks} (data fetching logic), and \texttt{src/types} (TypeScript DTOs matching backend entities). Keeping these separate means business logic and presentation concerns do not bleed into each other.

% ============================================================
\section{User Role-Based Access Control (RBAC)}
\label{sec:routing-access}

The frontend uses Role-Based Access Control (RBAC) to separate public and protected routes.

\subsection{Role Partitioning}
The application supports four user profiles: \textbf{Farmers} (verification and dispute submission), \textbf{Agricultural Officers} (inspection queues and evidence review), \textbf{Administrators} (system governance and telemetry) and \textbf{Public Users} (anonymous checks).

\subsection{Route Protection Middleware}
A \texttt{RoleGuard} component checks JWT claims against required permissions before mounting any protected React component, preventing privilege escalation at the route level. The core implementation logic is detailed in Listing~\ref{lst:role_guard}.

\begin{lstlisting}[
    style=mystyle,
    language=TypeScript,
    caption={Core Logic of the Role-Based Route Protection Guard},
    label={lst:role_guard}
]
export const RoleGuard: React.FC<{ children: React.ReactNode; allowedRoles: Role[] }> = ({ children, allowedRoles }) => {
    const { token, role, isLoading } = useAuthStore();
    const router = useRouter();

    useEffect(() => {
        if (!isLoading) {
            if (!token) return router.replace('/login');
            if (role && !allowedRoles.includes(role)) return router.replace('/unauthorized');
        }
    }, [token, role, isLoading, router, allowedRoles]);

    if (isLoading || !token) return <Loader />;
    return <>{children}</>;
};
\end{lstlisting}

% ============================================================
\section{Authentication and Secure Persistence}
\label{sec:authentication}

Sessions are secured using a dual-token JWT exchange designed to limit exposure to XSS attacks \cite{masue2024quantifying}.

\subsection{Token Storage Architecture}
Short-lived access tokens are held in the Zustand in-memory store, keeping them out of reach of malicious scripts. Long-lived refresh tokens are persisted in encrypted storage or \texttt{HttpOnly} cookies, allowing background session renewal without exposing token values to client-side code.

\subsection{Automated Token Rotation}
Axios interceptors handle token expiry without user intervention. On a \texttt{401 Unauthorized} response, the interceptor pauses pending requests, runs a refresh cycle, and replays the original request with the new token. From the user's side, the session continues uninterrupted. The automated refresh interceptor configuration is shown in Listing~\ref{lst:axios_interceptor}.

\begin{lstlisting}[
    style=mystyle,
    language=TypeScript,
    caption={Axios Interceptor for Automated JWT Token Rotation},
    label={lst:axios_interceptor}
]
apiClient.interceptors.response.use(res => res, async (error) => {
    const origReq = error.config;
    if (error.response?.status === 401 && !origReq._retry) {
        origReq._retry = true;
        try {
            const res = await axios.post('/api/proxy/auth/refresh', { token: useAuthStore.getState().refreshToken });
            useAuthStore.getState().setTokens(res.data.accessToken, res.data.refreshToken);
            origReq.headers['Authorization'] = `Bearer ${res.data.accessToken}`;
            return apiClient(origReq);
        } catch (err) {
            useAuthStore.getState().logout();
            return Promise.reject(err);
        }
    }
    return Promise.reject(error);
});
\end{lstlisting}

% ============================================================
\section{API Integration and Type Safety}
\label{sec:api-integration}

TypeScript DTOs enforce compile-time verification of HTTP payloads across all integration modules Auth, Verification, Analytics, and Complaint Services. Each mirrors the corresponding ASP.NET Core schema, so data arriving at the frontend matches what the UI expects without runtime coercion or defensive casting.

% ============================================================
\section{State Management}
\label{sec:state-management}

State is split by volatility: server-driven data and client-local data are managed separately.

\subsection{Asynchronous Server State (TanStack Query)}
Network-driven data (dispute queues, telemetry) is managed by TanStack Query, which handles stale-while-revalidate caching, background polling, and optimistic UI updates. On slow connections, this keeps the interface usable even when the server is slow to respond \cite{tanstack2024}.

\subsection{Synchronous Client State (Zustand)}
UI state that doesn't touch the network (active theme, session context) lives in Zustand. Its atomic slices limit unnecessary re-renders, which matters on lower-end devices where battery and memory are constrained.

% ============================================================
\section{Responsive UI Design and Field Usability}
\label{sec:responsive-ui}

The UI is built mobile-first using Tailwind CSS. Most farmers will access the application on a mid-range Android handset, often outdoors, so the interface is designed around that constraint rather than treating it as an edge case \cite{singh2026impact}.

\subsection{Standardized UI Library}
The interface is built from a set of reusable primitives (Table~\ref{tab:ui-components}) to keep visual behaviour consistent across screens.

\begin{table}[htbp]
\centering
\caption{Standardized UI Components}
\label{tab:ui-components}
\small
\renewcommand{\arraystretch}{1.1}
\begin{tabularx}{\textwidth}{|p{3.8cm}|X|}
\hline
\textbf{Component} & \textbf{Role} \\
\hline
\textbf{Data Table} & Scrollable grids with sticky headers and dynamic pagination for inspection data. \\
\hline
\textbf{Status Badge} & Colour-coded indicators (e.g., Green/Genuine) for fast visual parsing. \\
\hline
\textbf{Loading Skeleton} & Animated placeholders that prevent layout shifts during network delays. \\
\hline
\textbf{Input Modal} & Focused overlays that trap keyboard navigation for accessible data entry. \\
\hline
\end{tabularx}
\end{table}

\subsection{Accessibility Enhancements}
The UI meets WCAG AAA contrast ratios for readability in bright outdoor conditions, uses a minimum touch target size of $48\times 48$px to reduce input errors, and includes full ARIA semantic tagging for screen-reader compatibility.
